Multimodal Sentiment Analysis
\citep{zadeh2016mosi,zadeh2018mosei,liu2022simsv2}
predicts the sentiment intensity of an utterance from three loosely
aligned streams: language, acoustic prosody and facial dynamics.
The dominant design pattern is to use a strong text encoder
(BERT-style) as the primary backbone and inject acoustic / visual
context through cross-modal attention or low-rank fusion
\citep{rahman2020magbert,tsai2019mult,yu2021selfmm,han2021mmim}.
Recent work adds an information-bottleneck (IB) regulariser to compress
each modality into a task-relevant subspace
\citep{ithp2024,wu2023dbf,cmib2023,camib2026}, which improves the
state-of-the-art MAE on \textsc{CMU-MOSI} and \textsc{CMU-MOSEI}.

Three weaknesses persist across this line of work:
\textbf{(W1)} Cross-modal attention treats every auxiliary token as
equally informative.
A whispered \textit{``yeah''} and a screamed \textit{``yeah''} have
identical text content but very different acoustic information; current
key/value gating does not see this distinction.
\textbf{(W2)} The injection strength of auxiliary tokens into the
primary stream is fixed regardless of \emph{cross-modal consistency}.
A smiling face during a sarcastic verbal utterance contradicts the
language signal, but is added to the primary representation with the
same weight as a corroborating face.
\textbf{(W3)} Most architectures pre-commit to a fixed primary
modality, usually text.
Recent dynamic-routing approaches such as MODS \citep{wang2023cross}
introduce a learned modality selector, but the selector is trained from
the task gradient alone and tends to collapse onto the easy text-only
prior.

We argue that all three weaknesses share a single root cause: the
absence of a calibrated, sample-wise notion of \emph{information
quality} that can be propagated through the fusion pipeline.
We provide one.
The variational IB encoder used by ITHP and CyIN already produces a
per-token log-variance $\log\sigma^{2,m}$, and we observe that the
corresponding confidence
$\conf^{m}=\sigmoid(-\log\sigma^{2,m})$
is a natural, end-to-end-learnable measure of how task-relevant a token
is.
We use it once to modulate \emph{everything} downstream.

Concretely, our model \methodname{} (Figure~\ref{fig:overview}) uses
$\conf^{m}$ at three different scales:
\begin{itemize}
\itemsep0.1em
\item \textbf{Token level} -- IB-Guided Cross-Attention biases scores
  by $\gamma\log\bar c^{a}$ and gates values by $\bar c^{a}$, so noisy
  auxiliary positions are simultaneously down-weighted in attention and
  in contribution (\S\ref{sec:method:ibca}).
\item \textbf{Sample level} -- Alignment-Modulated Injection scales
  each cross-attention output by the cosine similarity of the primary
  and auxiliary bottleneck representations, lower-bounded at $0.3$;
  the injection therefore contracts when modalities disagree
  (\S\ref{sec:method:align}).
\item \textbf{Global level} -- a Gumbel-softmax modality selector
  picks one of the three modalities as the primary stream per sample,
  supervised in stage 2 by a routing-IB loss whose teacher signal is
  the per-modality task error in the bottleneck space
  (\S\ref{sec:method:mselector}).
\end{itemize}
This single mechanism resolves \textbf{W1--W3} simultaneously and is
fully differentiable through a single $\log\sigma^{2}$ pathway.

Our contributions are: (i)~the \emph{IB-confidence-as-modulation}
principle that unifies token, sample and global filtering through one
back-propagatable signal; (ii)~the IB-Guided Cross-Attention
mechanism with score bias and value gate; (iii)~the
Alignment-Modulated Injection rule for sample-level cross-modal
consistency control; (iv)~a stage-2 routing-IB loss that supervises
the modality selector with per-modality task error; and
(v)~state-of-the-art results on \textsc{CMU-MOSI} (MAE 0.5923) and
\textsc{CMU-MOSEI} (MAE 0.4941) under careful Optuna search, with a
competitive trade-off on \textsc{CH-SIMS}~v2.
The paper is organised as follows.
\S\ref{sec:related} reviews multimodal sentiment fusion and IB-based
methods.
\S\ref{sec:method} formalises \methodname{}.
\S\ref{sec:exp} reports the main results, ablations and
hyper-parameter search.
\S\ref{sec:analysis} analyses routing behaviour and the
MAE-vs-accuracy trade-off.
\S\ref{sec:conclusion} concludes; limitations and ethics statements
follow.
